\documentclass{article}
\usepackage{booktabs}
\usepackage{listings}
\usepackage{hyperref}

\lstset{
  language=Python,
  basicstyle=\ttfamily\small,
  breaklines=true,
  frame=single
}

\begin{document}

\section*{Model Card: Ztf Wd Catalog Monthly Gaia Warm source\_id Score Prediction Model}

\subsection*{Model Details}
\begin{itemize}
  \item \textbf{Model type:} Bayesian Hierarchical Regression with Time-Varying Effects
  \item \textbf{Version:} 0.23.0
  \item \textbf{Authors:} Ztf Wd Catalog Monthly Gaia Warm Prediction Pipeline
  \item \textbf{Created:} 2026-08-03
  \item \textbf{Last updated:} 2026-08-03
\end{itemize}

\subsection*{Intended Use}
This model is intended for:

- Research on source id score trajectories
- Exploration of month id-to-month id score patterns
- Educational demonstration of Bayesian hierarchical modeling

\subsubsection*{Out-of-Scope Use}
This model should NOT be used for:

- High-stakes or operational decisions without human review
- Real-time prediction systems in production environments

\subsection*{Training Data}
\begin{itemize}
  \item \textbf{Dataset:} ztf\_wd\_catalog\_monthly\_gaia\_warm
  \item \textbf{Size:} 49,422 albums
  \item \textbf{Description:} Sequential month id-level records grouped by source id from the 'ztf\_wd\_catalog\_monthly\_gaia\_warm' dataset, with the target score (mag\_binned) bounded to [10, 20].
  \item \textbf{Preprocessing:} Leak-safe within-source id temporal splitting with source id-disjoint secondary checks, minimum observation-count filtering, features standardized to zero mean and unit variance.
\end{itemize}

\subsection*{Model Architecture}
Grouping entity: source\_id; sequential event: month\_id. Internal parameter names use 'artist' for the grouping entity.

Bayesian hierarchical regression with four key components:

1. **Hierarchical artist effects**: Partial pooling across artists for robust estimation of artist quality. Non-centered parameterization via LocScaleReparam avoids funnel geometry.

2. **Time-varying slopes**: Artist quality modeled as a random walk, allowing career trajectories to evolve over time.

3. **AR(1) structure**: Album-to-album dependencies captured via autoregressive term, modeling momentum effects where consecutive albums tend to have correlated scores.

4. **Heteroscedastic observation noise** (sigma\_ref parameterization): Albums with more reviews have lower observation noise. The model samples sigma\_ref (noise at the median review count n\_ref) and derives per-observation noise as: sigma\_obs = sigma\_ref * n\_ref\textasciicircum{}n\_exponent, then sigma\_i = sigma\_obs / n\_reviews\_i\textasciicircum{}n\_exponent. This reparameterization breaks the multiplicative funnel between sigma\_obs and n\_exponent that causes divergent transitions in MCMC sampling.

Mathematical form:
- y\_ij \textasciitilde{} StudentT(df=4, mu\_ij, sigma\_i) (likelihood family configurable via --likelihood-family)
- mu\_ij = artist\_effect\_jt + X\_ij @ beta + rho * (prev\_score\_ij - ar\_center)
- artist\_effect\_jt evolves via random walk from initial effect
- sigma\_i = sigma\_obs / n\_reviews\_i\textasciicircum{}n\_exponent (heteroscedastic mode)

Since 0.5.0 the likelihood operates on the offset-logit transformed target by default (--target-transform offset\_logit), and since 0.6.0 the default feature block includes a stacked-GBM offset (gbm\_offset) with genre-level pooling where the dataset provides a group column. The per-run values are recorded in the hyperparameters table.

\subsubsection*{Prior Distributions}
Prior distributions are weakly informative, chosen to regularize inference while allowing the data to dominate:

- **mu\_artist** \textasciitilde{} Normal(0.7227843999862671, 1.0): Centered at 0.7227843999862671 because artist effects represent deviations from feature-based predictions. Scale of 1.0 permits the population center to shift by \textasciitilde{}1.0 SD on the standardized score scale.
- **sigma\_artist** \textasciitilde{} HalfNormal(0.5): Scale of 0.5 encourages moderate partial pooling. Implies most artist effects within +/-1.0, consistent with observed between-artist spread.
- **sigma\_rw** \textasciitilde{} HalfNormal(0.1): Scale of 0.1 produces smooth career trajectories where album-to-album quality changes are small relative to overall artist variation.
- **rho** \textasciitilde{} TruncatedNormal(0.0, 0.3, -0.99, 0.99): Centered at 0.0 with scale 0.3, allowing moderate autoregressive momentum without strong prior commitment to direction.
- **beta** \textasciitilde{} Normal(0.0, 1.0): Scale of 1.0 is weakly informative for standardized features, allowing data to determine effect sizes.
- **sigma\_obs** \textasciitilde{} HalfNormal(0.5): Scale of 0.5 allows data to determine observation-level noise.


**Prior Predictive Check**: Prior predictive simulation (n\_samples=500) shows 74.2\% of prior-implied predictions fall within [10.0, 20.0]. Summary: mean=15.9, sd=3.8, range=[9.5, 20.5].

\subsubsection*{Hyperparameters}
\begin{tabular}{ll}
\toprule
Parameter & Value \\
\midrule
mu\_artist\_loc & 0.0 \\
mu\_artist\_scale & 1.0 \\
sigma\_artist\_scale & 0.5 \\
sigma\_rw\_scale & 0.1 \\
rho\_loc & 0.0 \\
rho\_scale & 0.3 \\
beta\_loc & 0.0 \\
beta\_scale & 1.0 \\
sigma\_obs\_scale & 1.0 \\
sigma\_ref\_scale & 1.0 \\
n\_exponent\_default & 0.0 \\
n\_features & 8 \\
n\_artists & 928 \\
max\_albums & 100 \\
num\_chains & 4 \\
num\_warmup & 3000 \\
num\_samples & 3000 \\
chain\_method & sequential \\
target\_accept\_prob & 0.9 \\
max\_tree\_depth & 10 \\
target\_transform & offset\_logit \\
likelihood\_family & studentt \\
entity\_group\_pooling & False \\
gbm\_offset & False \\
\bottomrule
\end{tabular}

\subsection*{Evaluation Results}

\subsubsection*{Convergence Diagnostics}
Convergence status: PASSED

- R-hat (max): 1.0000 (threshold: < 1.01)
- ESS bulk (min): 5,100
- ESS tail (min): 6,603
- Divergent transitions: 0

\subsubsection*{Calibration}
Credible interval coverage:
- 80\% CI: 77.9\% empirical coverage, mean width=0.07
- 95\% CI: 92.3\% empirical coverage, mean width=0.13

**Posterior Predictive Checks:**
- mean: T(y\_obs)=16.91, p=0.930 (MC SE: 0.002)
- sd: T(y\_obs)=0.89, p=0.197 (MC SE: 0.004)
- skewness: T(y\_obs)=-1.27, p=0.722 (MC SE: 0.004)
- min: T(y\_obs)=12.81, p=0.442 (MC SE: 0.005)
- max: T(y\_obs)=18.94, p=0.001 (MC SE: 0.000)
- q10: T(y\_obs)=15.64, p=0.855 (MC SE: 0.003)
- q50: T(y\_obs)=17.12, p=0.864 (MC SE: 0.003)
- q90: T(y\_obs)=17.82, p=0.941 (MC SE: 0.002)

\subsubsection*{Predictive Performance}
Point prediction metrics:

- MAE: 0.02
- RMSE: 0.04
- R-squared: 0.998
- Held-out ELPD (test lppd): 1847.4 (SE: 36.2)

Held-out ELPD (test lppd): 1847.4

\subsection*{Limitations}
\begin{itemize}
  \item Dynamic source id trajectories are learned only when an source id has at least 2 training month ids.
  \item Score predictions are probabilistic and should not be treated as ground truth.
  \item Review the AOTY model card limitations for the statistical caveats of the shared model architecture (bounded target vs. symmetric likelihood, convergence budget).
\end{itemize}

\subsection*{Ethical Considerations}
\begin{itemize}
  \item Predictions are for research and exploration, not operational decisions.
  \item Historical biases in the recorded mag\_binned values will be reflected in predictions.
\end{itemize}

\subsection*{How to Use}

\subsubsection*{Loading the Model}
\begin{lstlisting}
from pathlib import Path

from panelcast.models.bayes.io import load_manifest, load_model

# Load the current ztfgcgw-score model referenced by models/manifest.json
manifest = load_manifest(Path("models"))
model_name = manifest.current["ztfgcgw_score"]
idata = load_model(Path("models") / model_name)
\end{lstlisting}

\subsubsection*{Making Predictions}
\begin{lstlisting}
from panelcast.models.bayes.predict import (
    extract_posterior_samples,
    predict_new_entity,
)
import jax.numpy as jnp

posterior_samples = extract_posterior_samples(idata)
n_features = int(posterior_samples["ztfgcgw_beta"].shape[-1])
X_new = jnp.zeros((1, n_features), dtype=jnp.float32)
pred = predict_new_entity(
    posterior_samples=posterior_samples,
    X_new=X_new,
    prev_score=jnp.array([15], dtype=jnp.float32),
    n_reviews_new=jnp.array([100.0], dtype=jnp.float32),
    prefix="ztfgcgw_",
    target_bounds=(10, 20),
)
\end{lstlisting}

\subsubsection*{Interpreting Results}
\begin{lstlisting}
import numpy as np

# Extract prediction statistics from posterior predictive draws
y_samples = np.asarray(pred['y']).ravel()
pred_mean = float(np.mean(y_samples))
pred_std = float(np.std(y_samples))
ci_95 = np.percentile(y_samples, [2.5, 97.5])

print(f"Predicted score: {pred_mean:.1f} +/- {pred_std:.1f}")
print(f"95% CI: [{ci_95[0]:.1f}, {ci_95[1]:.1f}]")
\end{lstlisting}

\end{document}